\chapter{Pecahan: Langkah Pertama}\label{ch:g6:fractions}

Membagi tiga pizza secara adil kepada empat orang memberi setiap
orang\,\dots\ bukan bilangan bulat pizza. \omterm{def:g6:fractions:def}{Pecahan} adalah bilangan yang
ditemukan untuk membagi. Bab ini membangun gambarannya: \omterm{def:g6:fractions:def}{pecahan}
sebagai bagian dari suatu keseluruhan, sebagai titik pada garis
bilangan, dan sebagai hasil pembagian yang tepat.

\section{Apa arti sebuah pecahan}

\begin{definition}[Pecahan]\label{def:g6:fractions:def}
Potonglah satu satuan menjadi $b$ bagian yang sama, lalu ambil $a$ di
antaranya: banyaknya yang diperoleh adalah \emph{pecahan}\index{pecahan}
\[
\frac{a}{b}
\qquad
\begin{array}{l}
a \text{: \emph{pembilang} --- berapa bagian yang kita ambil;}\\
b \text{: \emph{penyebut} --- satuannya dipotong menjadi berapa bagian.}
\end{array}
\]
\end{definition}

\begin{omfigure}
\begin{tikzpicture}[scale=0.55]
  \foreach \i in {0,...,3} \draw[omExo] (\i*2,0) rectangle (\i*2+2,1.2);
  \foreach \i in {0,...,2} \fill[omDef!40] (\i*2+0.05,0.05)
    rectangle (\i*2+1.95,1.15);
  \draw[thick] (0,0) rectangle (8,1.2);
  \node[below, font=\small] at (4,-0.15) {$\frac34$ batang itu: potong
  jadi $4$, ambil $3$};
  \begin{scope}[xshift=11.5cm]
  \foreach \i in {0,...,7} \draw[omExo] (\i,0) rectangle (\i+1,1.2);
  \foreach \i in {0,...,4} \fill[omThm!35] (\i+0.05,0.05)
    rectangle (\i+0.95,1.15);
  \draw[thick] (0,0) rectangle (8,1.2);
  \node[below, font=\small] at (4,-0.15) {$\frac58$ batang itu: potong
  jadi $8$, ambil $5$};
  \end{scope}
\end{tikzpicture}

{\small Dua \omterm{def:g6:fractions:def}{pecahan} dari batang yang sama. \omterm{def:g4:fractions:def}{Penyebut} memberi tahu
sehalus apa potongannya, \omterm{def:g4:fractions:def}{pembilang} memberi tahu berapa yang diambil.}
\end{omfigure}

\begin{example}\label{ex:g6:fractions:morethanone}
Sebuah \omterm{def:g6:fractions:def}{pecahan} dapat lebih besar daripada $1$: $\frac74$ berarti tujuh
\omterm{def:g3:division:half}{seperempat} --- satu satuan utuh (empat \omterm{def:g3:division:half}{seperempat}) dan tiga \omterm{def:g3:division:half}{seperempat} lagi:
\[
\frac74 = 1 + \frac34 .
\]
\end{example}

\begin{proposition}[Pecahan pada garis bilangan]\label{prop:g6:fractions:line}
Untuk meletakkan $\frac ab$ pada garis bilangan, potonglah setiap
selang satuan menjadi $b$ bagian yang sama lalu bilanglah $a$ bagian dari $0$.
\end{proposition}

\begin{omfigure}
\begin{tikzpicture}[scale=1.5]
  \draw[-Stealth] (-0.2,0) -- (7.6,0);
  \foreach \i in {0,4,8,12,16,20,24,28}
    \draw ({\i*0.25},-0.09) -- ({\i*0.25},0.09);
  \foreach \i in {0,...,28} \draw ({\i*0.25},-0.05) -- ({\i*0.25},0.05);
  \node[below, font=\small] at (0,-0.1) {$0$};
  \node[below, font=\small] at (1,-0.1) {$1$};
  \node[below, font=\small] at (2,-0.1) {$2$};
  \node[below, font=\small] at (3,-0.1) {$3$};
  \node[below, font=\small] at (4,-0.1) {$4$};
  \node[below, font=\small] at (5,-0.1) {$5$};
  \node[below, font=\small] at (6,-0.1) {$6$};
  \node[below, font=\small] at (7,-0.1) {$7$};
  \fill[omDef] (0.75,0) circle (1.6pt)
    node[above, font=\small] {$\frac34$};
  \fill[omThm] (1.75,0) circle (1.6pt)
    node[above, font=\small] {$\frac74$};
  \fill[omMeth] (5.5,0) circle (1.6pt)
    node[above, font=\small] {$\frac{22}{4}$};
\end{tikzpicture}

{\small Garis yang berskala seperempatan: $\frac34$ ada sebelum $1$,
$\frac74$ ada di antara $1$ dan $2$, dan $\frac{22}{4} = 5 + \frac24$
ada tepat di tengah antara $5$ dan $6$.}
\end{omfigure}

\begin{theorem}[Pecahan adalah sebuah pembagian]\label{thm:g6:fractions:quotient}
\omterm{def:g6:fractions:def}{Pecahan} $\frac ab$ adalah bilangan yang, jika dikalikan $b$, memberi
$a$:
\[
\frac ab \times b = a .
\]
Dengan kata lain, $\frac ab$ adalah \omterm{def:g3:division:remainder}{hasil bagi} $a$ oleh $b$ yang tepat
--- bahkan ketika pembagiannya ``tidak pas''.
\end{theorem}
\admitted

\begin{example}\label{ex:g6:fractions:pizzas}
Tiga pizza untuk empat orang: setiap orang mendapat $\frac34$ pizza.
Periksa dengan teoremanya: $4$ bagian sebesar $\frac34$ menjadi
$\frac34 \times 4 = 3$ pizza. Sebagian \omterm{def:g6:fractions:def}{pecahan} adalah \omterm{def:g6:decimals:places}{bilangan desimal}
($\frac34 = 0.75$), sebagian lagi tidak: $\frac13 = 0.333\dots$ tidak
pernah berhenti. \omterm{def:g6:fractions:def}{Pecahan} itulah nilai yang \emph{tepat}; $0.33$
hanyalah hampiran.
\end{example}

\section{Pecahan yang senilai}

\begin{proposition}[Pecahan senilai]\label{prop:g6:fractions:equal}
Sebuah \omterm{def:g6:fractions:def}{pecahan} tidak berubah ketika \omterm{def:g4:fractions:def}{pembilang} dan \omterm{def:g4:fractions:def}{penyebutnya}
sama-sama dikalikan (atau dibagi) dengan bilangan bukan \omterm{def:g1:counting:zero}{nol} yang sama:
\[
\frac{a}{b} = \frac{a \times k}{b \times k} .
\]
\end{proposition}

\begin{proof}[Gagasan pembuktian]
Memotong setiap bagian menjadi $k$ potongan yang lebih kecil
mengalikan banyaknya potongan yang diambil maupun banyaknya potongan
seluruhnya dengan $k$, tetapi banyaknya yang diambil tetap sama.
Gambarlah yang paling jelas menjelaskannya:
\end{proof}

\begin{omfigure}
\begin{tikzpicture}[scale=0.6]
  \foreach \i in {0,1,2} \draw[omExo] (\i*2.4,0) rectangle (\i*2.4+2.4,1.2);
  \foreach \i in {0,1} \fill[omDef!40] (\i*2.4+0.05,0.05)
    rectangle (\i*2.4+2.35,1.15);
  \draw[thick] (0,0) rectangle (7.2,1.2);
  \node[below, font=\small] at (3.6,-0.15) {$\frac23$};
  \begin{scope}[xshift=10.5cm]
  \foreach \i in {0,...,5} \draw[omExo] (\i*1.2,0) rectangle (\i*1.2+1.2,1.2);
  \foreach \i in {0,...,3} \fill[omDef!40] (\i*1.2+0.05,0.05)
    rectangle (\i*1.2+1.15,1.15);
  \draw[thick] (0,0) rectangle (7.2,1.2);
  \node[below, font=\small] at (3.6,-0.15) {$\frac46$};
  \end{scope}
\end{tikzpicture}

{\small $\frac23 = \frac46$: memotong setiap sepertiga menjadi dua
melipatduakan bilangan di atas dan di bawah, tanpa mengubah bagian yang terwarnai.}
\end{omfigure}

\begin{method}[Menyederhanakan pecahan]\label{met:g6:fractions:simplify}
Carilah bilangan yang membagi habis \omterm{def:g4:fractions:def}{pembilang} \emph{dan} \omterm{def:g4:fractions:def}{penyebutnya},
lalu bagilah keduanya. Ulangi sampai tidak ada \omterm{def:g5:division:multiple}{pembagi} bersama lagi:
\[
\frac{12}{18} = \frac{12 \div 2}{18 \div 2} = \frac{6}{9}
= \frac{6 \div 3}{9 \div 3} = \frac{2}{3}.
\]
\end{method}

\section{Mengambil pecahan dari suatu banyaknya}

\begin{method}[Pecahan dari suatu banyaknya]\label{met:g6:fractions:ofquantity}
Untuk menghitung $\frac ab$ dari suatu banyaknya, bagilah banyaknya itu
dengan $b$ (satu bagian), lalu kalikan dengan $a$ (banyaknya bagian):
\[
\frac ab \text{ dari } Q = (Q \div b) \times a .
\]
(Mengalikan dulu dengan $a$ lalu membagi dengan $b$ memberi hasil yang
sama --- pilihlah urutan yang membuat bilangannya lebih enak.)
\end{method}

\begin{example}\label{ex:g6:fractions:ofquantity}
Hitunglah $\frac35$ dari $40$ euro, selangkah demi selangkah:
\begin{enumerate}
  \item seperlima dari $40$: $40 \div 5 = 8$;
  \item tiga perlima: $8 \times 3 = 24$.
\end{enumerate}
Jadi $\frac35$ dari $40$ euro adalah $24$ euro. Periksa: \omterm{def:g3:division:remainder}{sisanya}
$\frac25$ adalah $16$ euro, dan $24 + 16 = 40$.
\end{example}

\section{Latihan}

\begin{exercise}[$\star$]\label{exo:g6:fractions:1}
Untuk setiap gambaran berikut, tulislah \omterm{def:g6:fractions:def}{pecahannya}: sebuah kue
dipotong menjadi $8$ iris, $3$ dimakan; sebatang cokelat berisi $12$
kotak, $7$ dimakan; sebuah pai dipotong menjadi $6$, dan semua $6$ irisnya dimakan.
\end{exercise}

\begin{exercise}[$\star$]\label{exo:g6:fractions:2}
Gambarlah sebuah batang yang dipotong menjadi $5$ bagian sama besar
lalu warnai $\frac45$ bagiannya. Apakah $\frac45$ lebih kecil atau lebih besar daripada $1$? Dan $\frac65$?
\end{exercise}

\begin{exercise}[$\star$]\label{exo:g6:fractions:3}
Letakkan pada garis bilangan berskala sepertigaan: $\frac13$; \
$\frac53$; \ $2$; \ $\frac73$. Dua bilangan mana di antaranya yang
berjarak sama dari $2$?
\end{exercise}

\begin{exercise}[$\star$]\label{exo:g6:fractions:4}
Tulislah setiap \omterm{def:g6:fractions:def}{pecahan} sebagai bilangan bulat ditambah \omterm{def:g6:fractions:def}{pecahan} yang
lebih kecil daripada $1$ (seperti pada $\frac74 = 1 + \frac34$):
\[
\frac{9}{4}, \qquad \frac{13}{5}, \qquad \frac{12}{6}, \qquad
\frac{25}{8} .
\]
\end{exercise}

\begin{exercise}[$\star$]\label{exo:g6:fractions:5}
Lengkapilah agar \omterm{def:g6:fractions:def}{pecahannya} senilai:
\[
\frac{2}{3} = \frac{?}{12}, \qquad
\frac{15}{20} = \frac{3}{?}, \qquad
\frac{7}{10} = \frac{21}{?}, \qquad
\frac{?}{6} = \frac{10}{12}.
\]
\end{exercise}

\begin{exercise}[$\star$]\label{exo:g6:fractions:6}
Sederhanakan sedapat mungkin: $\dfrac{6}{8}$; \quad $\dfrac{15}{25}$;
\quad $\dfrac{18}{24}$; \quad $\dfrac{35}{7}$.
\end{exercise}

\begin{exercise}[$\star$]\label{exo:g6:fractions:7}
Hitunglah: $\frac12$ dari $86$; \quad $\frac34$ dari $60$; \quad
$\frac25$ dari $35$; \quad $\frac{7}{10}$ dari $250$.
\end{exercise}

\begin{exercise}[$\star$]\label{exo:g6:fractions:8}
Mana yang lebih besar: $\frac12$ atau $\frac13$ (dari kue yang sama)?
$\frac25$ atau $\frac35$? $\frac34$ atau $\frac78$ (tulislah keduanya
dengan \omterm{def:g4:fractions:def}{penyebut} $8$)?
\end{exercise}

\begin{exercise}[$\star\star$]\label{exo:g6:fractions:9}
Sebutkan nilai desimal \omterm{def:g6:fractions:def}{pecahan} yang punya nilai desimal, lalu katakan
\omterm{def:g6:fractions:def}{pecahan} mana yang tidak punya: $\frac12$; \ $\frac34$; \ $\frac13$; \
$\frac{7}{10}$; \ $\frac{9}{4}$.
\end{exercise}

\begin{exercise}[$\star\star$]\label{exo:g6:fractions:10}
Sebuah kelas berisi $28$ murid. Tiga perempat di antaranya datang ke
sekolah dengan bus, dan sisanya berjalan kaki. Berapa murid yang
berjalan kaki? Kerjakan selangkah demi selangkah, lalu periksa bahwa
kedua kelompokmu berjumlah $28$.
\end{exercise}

\begin{exercise}[$\star\star$]\label{exo:g6:fractions:11}
Leo membelanjakan $\frac23$ tabungannya untuk permainan seharga $18$
euro. Berapa uang yang dimiliki Leo sebelumnya? (Sepertiga tabungannya
adalah setengah dari $18$\,\dots\ pikirkan mengapa, atau carilah dulu
besarnya satu bagian.)
\end{exercise}

\begin{exercise}[$\star\star\star$]\label{exo:g6:fractions:12}
Sebuah tangki terisi $\frac58$ kapasitasnya. Setelah dipakai $10$
liter, tangki itu terisi setengah. Berapa \omterm{def:g2:measure:units}{kapasitas} tangki itu?
(\omterm{def:g6:fractions:def}{Pecahan} tangki yang mana yang diwakili $10$ liter itu?)
\end{exercise}

\section{Soal: Batang cokelat yang tidak pernah habis}

\begin{problem}[{Soal akhir pekan --- setengah dari setengah: jumlah
$\frac12 + \frac14 + \frac18 + \dots$ merayap menuju $1$ tanpa pernah
menyentuhnya}]
\label{pb:g6:fractions:1}
Ambillah sebatang cokelat. Makanlah setengahnya. Lalu makanlah
setengah dari \emph{yang tersisa}. Lalu setengah lagi dari yang
tersisa, dan seterusnya. Dua hal seakan-akan benar sekaligus:
cokelatnya tidak pernah habis (selalu ada yang tersisa), tetapi kamu
memakan hampir seluruhnya. Soal ini mengubah kedua perasaan itu
menjadi pernyataan yang tepat tentang \omterm{def:g6:fractions:def}{pecahan} --- dan di tengah jalan
bertemu seorang pelari dari Yunani kuno serta teka-teki pembagian yang
diselesaikan hanya dengan pisau yang cuma bisa membelah dua.

\medskip
\noindent\textbf{Bagian I --- Setengah dari setengah.}
Batang cokelat itu berbentuk persegi panjang berisi $16$ kotak sama besar.
\begin{enumerate}
  \item Gigitan pertama: kamu memakan setengah batang itu. Berapa
        kotak? Tulislah gigitan itu sebagai \omterm{def:g6:fractions:def}{pecahan} dari batangnya
        dengan dua cara, dengan \omterm{def:g4:fractions:def}{penyebut} $2$ dan dengan \omterm{def:g4:fractions:def}{penyebut} $16$
        (\cref{prop:g6:fractions:equal}).
  \item Gigitan kedua: setengah dari yang tersisa. Berapa kotak?
        Tunjukkan, dengan gambarnya atau dengan \omterm{def:g6:fractions:def}{pecahan} senilai, bahwa
        gigitan itu adalah $\frac14$ seluruh batang: \emph{setengah
        dari setengah adalah \omterm{def:g3:division:half}{seperempat}}.
  \item Gigitan ketiga dan keempat, dengan aturan yang sama: sebutkan
        masing-masing dalam kotak dan sebagai \omterm{def:g6:fractions:def}{pecahan} dari batangnya.
        Apa yang terjadi pada \omterm{def:g4:fractions:def}{penyebutnya} dari satu gigitan ke
        gigitan berikutnya?
  \item Setelah keempat gigitan itu, berapa kotak yang sudah dimakan
        seluruhnya? Berapa \omterm{def:g6:fractions:def}{pecahan} batangnya, dan berapa \omterm{def:g6:fractions:def}{pecahan} yang
        tersisa?
  \item Pertanyaan 4 mengatakan, ditulis dalam perenambelasan:
        \[
        \frac{8}{16} + \frac{4}{16} + \frac{2}{16} + \frac{1}{16}
        = \frac{15}{16},
        \qquad\text{yaitu}\qquad
        \frac12 + \frac14 + \frac18 + \frac{1}{16}
        = \frac{15}{16} .
        \]
        Periksalah keempat penulisan ulang itu
        ($\frac12 = \frac{8}{16}$, dan seterusnya).
\end{enumerate}

\medskip
\noindent\textbf{Bagian II --- Tangga menuju $1$.}
Agar dapat terus menggigit, bayangkan batang yang lebih halus berisi $64$ kotak.
\begin{enumerate}[resume]
  \item Sebutkan besar keenam gigitan yang pertama, dalam kotak
        ($32$, lalu\,\dots), dan periksalah bahwa setelah gigitan
        keenam tersisa tepat satu kotak. Berapa \omterm{def:g6:fractions:def}{pecahan} batangnya yang
        tersisa?
  \item Setelah setiap gigitan, berapa \omterm{def:g6:fractions:def}{pecahan} batangnya yang tersisa?
        Tulislah daftarnya ($\frac12$, $\frac14$, \dots) sampai
        gigitan keenam, lalu jelaskan polanya. Kalau sebuah pisau
        khayalan dapat membelah dua bahkan kotak yang terakhir, apa
        yang akan tersisa setelah gigitan ketujuh?
  \item Jelaskan mengapa \omterm{def:g1:addition:def}{jumlah} yang dimakan \emph{tidak pernah} dapat
        mencapai $1$, sebanyak apa pun gigitan yang dilakukan.
        (Bandingkan dengan dua puluh angka sembilan pada
        \cref{pb:g6:decimals:1}: cerita yang sama?)
  \item Namun \omterm{def:g1:addition:def}{jumlahnya} melewati sasaran mana pun di bawah $1$. Seorang
        teman menantangmu memakan lebih dari $\frac{99}{100}$ batang
        itu. Mana yang lebih kecil, $\frac{1}{128}$ atau
        $\frac{1}{100}$ --- dan sesudah gigitan yang mana tantangan
        itu dimenangkan?
  \item Pada garis bilangan dari $0$ sampai $1$ yang berskala
        perenambelasan, letakkan \omterm{def:g1:addition:def}{jumlah} yang dimakan setelah setiap
        gigitan dari keempat gigitan pertama: $\frac12$, $\frac34$,
        $\frac78$, $\frac{15}{16}$. Di mana setiap \omterm{def:g1:addition:def}{jumlah} yang baru
        mendarat, dibandingkan dengan \omterm{def:g1:addition:def}{jumlah} sebelumnya dan dengan $1$?
\end{enumerate}

\medskip
\noindent\textbf{Bagian III --- Pelari Zeno, dan pisau yang cuma
bisa membelah dua.}
\begin{enumerate}[resume]
  \item Dua puluh empat abad yang lalu, filsuf Yunani Zeno
        menceritakan kisah ini: untuk mencapai sebuah tembok, seorang
        pelari harus menempuh dulu setengah jaraknya, lalu setengah
        dari \omterm{def:g3:division:remainder}{sisanya}, lalu setengah lagi dari \omterm{def:g3:division:remainder}{sisanya}\,\dots\
        ``sehingga pelari itu tidak pernah sampai ke tembok!'' Dengan
        memakai pertanyaan 8 dan 9, katakan apa yang \emph{benar}
        dalam kisah Zeno, dan di mana perangkapnya. (Benarkah pelari
        itu memerlukan waktu selamanya untuk menempuh semua langkah
        tadi?)
  \item Teka-teki pembagian: bagilah $15$ batang cokelat yang identik
        secara adil kepada $16$ anak, hanya dengan potongan menjadi
        dua sama besar (dari satu batang, dari setengah batang, dari
        \omterm{def:g3:division:half}{seperempat} batang\,\dots). Jelaskan cara memotongnya: berapa
        batang yang dipotong menjadi setengahan, berapa menjadi
        seperempatan, menjadi perdelapanan, menjadi perenambelasan,
        sehingga setiap anak menerima satu potong dari tiap ukuran.
  \item Periksalah pembagianmu: tulislah bagian setiap anak sebagai
        \omterm{def:g1:addition:def}{jumlah} \omterm{def:g6:fractions:def}{pecahan}, hitunglah dalam perenambelasan, lalu
        pastikanlah bahwa $16$ bagian itu bersama-sama menghabiskan
        tepat $15$ batang tadi.
  \item Di kelas sebelah, $7$ batang dibagi rata kepada $8$ anak. Anak
        yang mana yang menerima cokelat lebih banyak --- satu dari
        yang $16$, atau satu dari yang $8$? (Tulislah kedua bagiannya
        dengan \omterm{def:g4:fractions:def}{penyebut} $16$, \cref{exo:g6:fractions:8}.)
  \item Penutupnya: seorang teman menyatakan bahwa dengan gigitan yang
        \emph{tak hingga} banyaknya orang akan memakan \emph{tepat}
        seluruh batang itu. Gambarlah sebuah persegi lalu warnai
        gigitannya di dalamnya: setengah persegi, lalu \omterm{def:g3:division:half}{seperempat},
        lalu seperdelapan\,\dots\ Apa yang kamu amati pada pojok yang
        tidak terwarnai? Nyatakan dengan cermat dua fakta yang benar
        pada seluruh soal ini: apa yang tersisa setelah gigitan yang
        banyaknya berhingga, dan sekecil apa \omterm{def:g3:division:remainder}{sisa} itu jadinya. (Apa
        arti tepat dari ``\omterm{def:g1:addition:def}{jumlah} tak hingga'' adalah cerita untuk buku
        Sekolah Menengah.)

\end{enumerate}
\end{problem}
